% Sections 2 (Background) and 3 (Posterior Schrodinger Bridges), compact version.
% Target length: about one page at \ninept. Include with % Sections 2 (Background) and 3 (Posterior Schrodinger Bridges), compact version.
% Target length: about one page at \ninept. Include with % Sections 2 (Background) and 3 (Posterior Schrodinger Bridges), compact version.
% Target length: about one page at \ninept. Include with % Sections 2 (Background) and 3 (Posterior Schrodinger Bridges), compact version.
% Target length: about one page at \ninept. Include with \input{sections_compact}.

\section{Background}
\label{sec:background}

Consider a sequence of $T$ disjoint tasks with datasets
$\mathcal{D}_1,\ldots,\mathcal{D}_T$ of i.i.d.\ samples and a model with
parameters $\theta\in\mathbb{R}^P$ shared across tasks, where the data of a task
are no longer available once it has been learned. Let $q_0(\theta)=p(\theta)$ be
the prior and $q_t(\theta)=p(\theta\mid\mathcal{D}_{1:t})$ the posterior after
task $t$. If the datasets are conditionally independent given $\theta$,
\begin{equation}
    q_{t+1}(\theta)\propto p(\mathcal{D}_{t+1}\mid\theta)\,q_t(\theta),
    \label{eq:sequential_bayes}
\end{equation}
so $q_t$ carries all the information from earlier tasks that the next update
requires. Since $q_t$ is intractable for neural networks, a replay-free learner
retains an approximation $\hat q_t$ instead. The update then defines a target
proportional to $p(\mathcal{D}_{t+1}\mid\theta)\,\hat q_t(\theta)$, which does not
directly provide samples, and errors in $\hat q_t$ propagate to all subsequent
updates~\cite{kessler2023sequential}. A replay-free learner must therefore both
represent the accumulated posterior and update this representation as new tasks
arrive.

\section{Posterior Schr\"odinger Bridges}
\label{sec:method}

PSB represents the posterior by an ensemble of networks, transports it to each
updated posterior by annealed Langevin dynamics, and fits a Schr\"odinger bridge
between consecutive ensembles that guides the next transport.

\subsection{Posterior representation and target}
\label{sec:representation}

We approximate the posterior after task $t$ by
\begin{equation}
    \hat q_t(\theta)=\sum_{k=1}^{K}w_t^{(k)}\,
    \mathcal{N}\!\left(\theta;\,\theta_t^{(k)},\,A_t^{-1}\right).
    \label{eq:mixture}
\end{equation}
Each term, which we call a component, is a Gaussian centered on a network
$\theta_t^{(k)}$, called a particle, with weight $w_t^{(k)}$. The components
cover distinct parameter solutions, and the shared diagonal precision
\begin{equation}
    A_t=\alpha I+\textstyle\sum_{j=1}^{t}\gamma^{\,t-j}N_jF_j
    \label{eq:precision}
\end{equation}
accumulates the Fisher information of each task, where $N_j=|\mathcal{D}_j|$ and
\begin{equation}
    F_j=\mathrm{diag}\Big(\mathbb{E}_{(x,y)\in\mathcal{D}_j}\mathbb{E}_{k}
    \big[\big(\nabla_\theta\log p(y\mid x,\theta_j^{(k)})\big)^{\odot2}\big]\Big)
    \label{eq:fisher}
\end{equation}
is the diagonal empirical Fisher of task $j$, with $\odot$ the elementwise
product. Here $\gamma\in(0,1]$ discounts older tasks and $\alpha>0$ is the
precision of a Gaussian prior. Since $N_jF_j$ is the Laplace precision of the
likelihood of task $j$, $A_t$ constrains each weight according to its importance
for earlier tasks. For the first task, the particles are He-uniform
initialized~\cite{he2015delving}, trained independently, and equally weighted,
treating the ensemble as approximate posterior
samples~\cite{lakshminarayanan2017ensembles,wilson2020bayesian}.

When task $t+1$ arrives, we introduce its likelihood through the tempered
posteriors
\begin{equation}
    \pi_{t,\beta}(\theta)\propto p(\mathcal{D}_{t+1}\mid\theta)^{\beta}\,
    \hat q_t(\theta),\qquad \beta\in[0,1].
    \label{eq:tempered}
\end{equation}
At $\beta=1$, this is Bayes' rule (\ref{eq:sequential_bayes}) applied to
$\hat q_t$, which defines the target $\tilde q_{t+1}=\pi_{t,1}$. Raising $\beta$
gradually avoids concentrating the weights on a few particles and moving the
particles over the full distance at once. The score of the tempered posteriors is
\begin{equation}
\begin{split}
    \nabla_\theta\log\pi_{t,\beta}(\theta)
    &= \beta\,\nabla_\theta\log p(\mathcal{D}_{t+1}\mid\theta)\\
    &\quad -A_t\Big(\theta-\textstyle\sum_{k=1}^{K}r_k(\theta)\,\theta_t^{(k)}\Big),
\end{split}
\label{eq:pathscore}
\end{equation}
where $r_k(\theta)\propto w_t^{(k)}\mathcal{N}(\theta;\theta_t^{(k)},A_t^{-1})$
is the probability that $\theta$ belongs to component $k$. The likelihood score
drives the particles toward the new task, and the prior score pulls each particle
toward the carried solutions it most likely belongs to, with a strength set by
$A_t$.

For a single component, the log-target is
$\log p(\mathcal{D}_{t+1}\mid\theta)-\tfrac12\|\theta-\theta_t\|_{A_t}^2$ up to a
constant, with $\|u\|_M^2=u^{\top}Mu$. Substituting (\ref{eq:precision}),
negating and dividing by $N_{t+1}$, maximizing the log-target is equivalent to
minimizing
\begin{equation}
    \bar\ell_{t+1}(\theta)
    +\tfrac{1}{2}\textstyle\sum_{j=1}^{t}\gamma^{\,t-j}\tfrac{N_j}{N_{t+1}}
    \|\theta-\theta_t\|_{F_j}^{2},
    \label{eq:ewc}
\end{equation}
where $\bar\ell_{t+1}$ is the mean negative log-likelihood of the new task and
the small isotropic term $\tfrac{\alpha}{2N_{t+1}}\|\theta-\theta_t\|^2$ is
omitted. This is the online EWC
objective~\cite{kirkpatrick2017ewc,schwarz2018progress} with penalty strength
fixed at $N_j/N_{t+1}$, so the EWC solution is the most probable parameter under
the target of PSB with a single particle.

\subsection{Annealed Langevin transport}
\label{sec:langevin}

As $\beta$ increases from $\beta_{m-1}$ to $\beta_m$, the particle weights are
multiplied by the change in the tempered likelihood,
\begin{equation}
    w^{(k)}\leftarrow w^{(k)}\,
    p\big(\mathcal{D}_{t+1}\mid\theta^{(k)}\big)^{\beta_m-\beta_{m-1}},
    \label{eq:reweight}
\end{equation}
with each increment as large as possible while the effective number of
particles $1/\sum_k(\bar w^{(k)})^2$ of the normalized weights stays above a
fixed fraction of $K$. When the weights become uneven, low-weight particles are
replaced by copies of high-weight ones and the weights are reset. The particles
then take preconditioned Langevin steps toward $\pi_{t,\beta_m}$,
\begin{equation}
    \theta'=\theta+\eta\,A_t^{-1}\nabla_\theta\log\pi_{t,\beta_m}(\theta)
    +\sqrt{2\eta}\,A_t^{-1/2}\xi,
    \label{eq:mala}
\end{equation}
with $\xi\sim\mathcal{N}(0,I)$ and an adapted step size $\eta$, each accepted
with a probability computed from the full task likelihood. The drift of these
dynamics is the score (\ref{eq:pathscore}), and the preconditioner makes the
steps cautious in directions constrained by earlier tasks. At $\beta=1$, the
particles form the terminal population $\{\Theta_1^{(k)}\}$.

\subsection{Fisher-metric Schr\"odinger bridge}
\label{sec:bridge}

Given two marginals and a reference process, the Schr\"odinger bridge (SB) is
the path measure closest to the reference in Kullback--Leibler (KL) divergence
that connects the marginals~\cite{leonard2014survey}. At each boundary, PSB
solves
\begin{equation}
    \mathbb{P}^\star_t=\mathop{\arg\min}_{\substack{\mathbb{P}:\ \mathbb{P}_0=\hat q_t\\
    \mathbb{P}_1=\tilde q_{t+1}}}\ \mathrm{KL}(\mathbb{P}\,\|\,\mathbb{Q}_t)
    \label{eq:psb}
\end{equation}
for the Fisher-metric Brownian reference
\begin{equation}
    \mathbb{Q}_t:\quad d\Theta_s=\sigma A_t^{-1/2}\,dW_s,\qquad s\in[0,1].
    \label{eq:fisherreference}
\end{equation}
For a process with added drift $u_s$, the KL divergence to $\mathbb{Q}_t$ is
$(2\sigma^2)^{-1}\mathbb{E}\!\int_0^1\|u_s\|_{A_t}^2ds$ plus a term fixed by the
initial marginal, so moving weights that earlier tasks constrain strongly is
costly. Solving (\ref{eq:psb}) at every boundary yields a chain of SBs
$\hat q_1\to\cdots\to\hat q_T$.

We solve (\ref{eq:psb}) by iterative Markovian fitting (IMF)~\cite{shi2023dsbm}.
Given a coupling of samples $(\Theta_0,\Theta_1)$ from $\hat q_t$ and the
terminal population, intermediate points are drawn from the conditional bridge
of $\mathbb{Q}_t$,
\begin{equation}
    \Theta_s=(1-s)\,\Theta_0+s\,\Theta_1+\sigma\sqrt{s(1-s)}\,A_t^{-1/2}\xi,
    \label{eq:conditionalbridge}
\end{equation}
and a forward drift is fitted by
\begin{equation}
    \min_\phi\ \mathbb{E}\left\|v^{\mathrm{f}}_\phi(\Theta_s,s)
    -\frac{\Theta_1-\Theta_s}{1-s}\right\|_{A_t}^{2},
    \label{eq:lossf}
\end{equation}
with $s$ uniform on $[0,1]$. A backward drift $v^{\mathrm{b}}_\psi$ is fitted
analogously with target $(\Theta_0-\Theta_s)/s$. New couplings are generated by alternately
simulating the fitted forward and backward processes, and since learned drifts
only approximate the IMF projections, we report the energy distance between
simulated and terminal populations.

A drift over all weights at once would be too large to fit from $K$ particles,
so the drifts act on neuron tokens. The token of unit $i$ in a fully connected
layer $\ell$ is its incoming weights and bias,
$[W^{(\ell)}_{i,:},\,b^{(\ell)}_i]$, and in a convolutional layer it is one
filter with its bias. Any vector $u$ over the weights is split in the same way,
with $[u]^{(\ell)}_i$ its token at unit $i$ of layer $\ell$. The forward drift is
predicted token by token,
\begin{equation}
    \big[v^{\mathrm{f}}_\phi(\Theta_s,s)\big]^{(\ell)}_i
    =f^{(\ell)}_\phi\!\left([\Theta_s]^{(\ell)}_i,\,s,\,[g_s]^{(\ell)}_i,\,
    [a_s]^{(\ell)}_i\right),
    \label{eq:tokendrift}
\end{equation}
where $g_s=\nabla_\theta\log p(\mathcal{D}_{t+1}\mid\Theta_s)$, $a_s$ is the
prior term of (\ref{eq:pathscore}) at $\Theta_s$, and $f^{(\ell)}_\phi$ is a
multilayer perceptron shared by the $n_\ell$ units of layer $\ell$ that receives
$s$ through Fourier features. Sharing turns each particle into $n_\ell$ training
examples. The score inputs tell the network how the new task and the carried
posterior pull on each unit, so the drift depends on this evidence rather than
on a task identifier and remains usable at the next boundary. A zero-initialized
output layer makes the untrained drift zero, so fitting starts from the
reference (\ref{eq:fisherreference}). The backward drift has the same form.

\subsection{Bridge-guided transitions}
\label{sec:transition}

The drift fitted at one boundary guides the transport at the next. The particles
are moved by the Euler--Maruyama discretization of the previous forward drift,
with $s_m=\beta_m$ and $\Delta s_m=s_m-s_{m-1}$,
\begin{equation}
    \Theta_{m}=\Theta_{m-1}+v^{\mathrm{f}}_\phi(\Theta_{m-1},s_{m-1})\,\Delta s_m
    +\sigma\sqrt{\Delta s_m}\,A_t^{-1/2}\xi,
    \label{eq:euler}
\end{equation}
and their weights are updated by
\begin{equation}
    w^{(k)}\leftarrow w^{(k)}\,
    \frac{\pi_{t,\beta_m}(\Theta_m)\,L_m(\Theta_{m-1}\mid\Theta_m)}
         {\pi_{t,\beta_{m-1}}(\Theta_{m-1})\,K_m(\Theta_m\mid\Theta_{m-1})},
    \label{eq:bridgeweight}
\end{equation}
where $K_m$ and $L_m$ are the Gaussian transition densities of (\ref{eq:euler})
and of the corresponding backward step~\cite{delmoral2006smc}. The terminal
population thus targets $\tilde q_{t+1}$ for any drift, while an accurate drift
keeps the weights nearly uniform and reduces the Langevin steps (\ref{eq:mala})
needed after each reset. The drift is zero at the first boundary and is refitted
by IMF after every transport. The terminal particles and the updated precision
(\ref{eq:precision}) form $\hat q_{t+1}$, predictions average the particles'
predictive distributions, and the stored state, $K$ particles, one diagonal
precision and the drift networks, does not grow with the number of tasks.
.

\section{Background}
\label{sec:background}

Consider a sequence of $T$ disjoint tasks with datasets
$\mathcal{D}_1,\ldots,\mathcal{D}_T$ of i.i.d.\ samples and a model with
parameters $\theta\in\mathbb{R}^P$ shared across tasks, where the data of a task
are no longer available once it has been learned. Let $q_0(\theta)=p(\theta)$ be
the prior and $q_t(\theta)=p(\theta\mid\mathcal{D}_{1:t})$ the posterior after
task $t$. If the datasets are conditionally independent given $\theta$,
\begin{equation}
    q_{t+1}(\theta)\propto p(\mathcal{D}_{t+1}\mid\theta)\,q_t(\theta),
    \label{eq:sequential_bayes}
\end{equation}
so $q_t$ carries all the information from earlier tasks that the next update
requires. Since $q_t$ is intractable for neural networks, a replay-free learner
retains an approximation $\hat q_t$ instead. The update then defines a target
proportional to $p(\mathcal{D}_{t+1}\mid\theta)\,\hat q_t(\theta)$, which does not
directly provide samples, and errors in $\hat q_t$ propagate to all subsequent
updates~\cite{kessler2023sequential}. A replay-free learner must therefore both
represent the accumulated posterior and update this representation as new tasks
arrive.

\section{Posterior Schr\"odinger Bridges}
\label{sec:method}

PSB represents the posterior by an ensemble of networks, transports it to each
updated posterior by annealed Langevin dynamics, and fits a Schr\"odinger bridge
between consecutive ensembles that guides the next transport.

\subsection{Posterior representation and target}
\label{sec:representation}

We approximate the posterior after task $t$ by
\begin{equation}
    \hat q_t(\theta)=\sum_{k=1}^{K}w_t^{(k)}\,
    \mathcal{N}\!\left(\theta;\,\theta_t^{(k)},\,A_t^{-1}\right).
    \label{eq:mixture}
\end{equation}
Each term, which we call a component, is a Gaussian centered on a network
$\theta_t^{(k)}$, called a particle, with weight $w_t^{(k)}$. The components
cover distinct parameter solutions, and the shared diagonal precision
\begin{equation}
    A_t=\alpha I+\textstyle\sum_{j=1}^{t}\gamma^{\,t-j}N_jF_j
    \label{eq:precision}
\end{equation}
accumulates the Fisher information of each task, where $N_j=|\mathcal{D}_j|$ and
\begin{equation}
    F_j=\mathrm{diag}\Big(\mathbb{E}_{(x,y)\in\mathcal{D}_j}\mathbb{E}_{k}
    \big[\big(\nabla_\theta\log p(y\mid x,\theta_j^{(k)})\big)^{\odot2}\big]\Big)
    \label{eq:fisher}
\end{equation}
is the diagonal empirical Fisher of task $j$, with $\odot$ the elementwise
product. Here $\gamma\in(0,1]$ discounts older tasks and $\alpha>0$ is the
precision of a Gaussian prior. Since $N_jF_j$ is the Laplace precision of the
likelihood of task $j$, $A_t$ constrains each weight according to its importance
for earlier tasks. For the first task, the particles are He-uniform
initialized~\cite{he2015delving}, trained independently, and equally weighted,
treating the ensemble as approximate posterior
samples~\cite{lakshminarayanan2017ensembles,wilson2020bayesian}.

When task $t+1$ arrives, we introduce its likelihood through the tempered
posteriors
\begin{equation}
    \pi_{t,\beta}(\theta)\propto p(\mathcal{D}_{t+1}\mid\theta)^{\beta}\,
    \hat q_t(\theta),\qquad \beta\in[0,1].
    \label{eq:tempered}
\end{equation}
At $\beta=1$, this is Bayes' rule (\ref{eq:sequential_bayes}) applied to
$\hat q_t$, which defines the target $\tilde q_{t+1}=\pi_{t,1}$. Raising $\beta$
gradually avoids concentrating the weights on a few particles and moving the
particles over the full distance at once. The score of the tempered posteriors is
\begin{equation}
\begin{split}
    \nabla_\theta\log\pi_{t,\beta}(\theta)
    &= \beta\,\nabla_\theta\log p(\mathcal{D}_{t+1}\mid\theta)\\
    &\quad -A_t\Big(\theta-\textstyle\sum_{k=1}^{K}r_k(\theta)\,\theta_t^{(k)}\Big),
\end{split}
\label{eq:pathscore}
\end{equation}
where $r_k(\theta)\propto w_t^{(k)}\mathcal{N}(\theta;\theta_t^{(k)},A_t^{-1})$
is the probability that $\theta$ belongs to component $k$. The likelihood score
drives the particles toward the new task, and the prior score pulls each particle
toward the carried solutions it most likely belongs to, with a strength set by
$A_t$.

For a single component, the log-target is
$\log p(\mathcal{D}_{t+1}\mid\theta)-\tfrac12\|\theta-\theta_t\|_{A_t}^2$ up to a
constant, with $\|u\|_M^2=u^{\top}Mu$. Substituting (\ref{eq:precision}),
negating and dividing by $N_{t+1}$, maximizing the log-target is equivalent to
minimizing
\begin{equation}
    \bar\ell_{t+1}(\theta)
    +\tfrac{1}{2}\textstyle\sum_{j=1}^{t}\gamma^{\,t-j}\tfrac{N_j}{N_{t+1}}
    \|\theta-\theta_t\|_{F_j}^{2},
    \label{eq:ewc}
\end{equation}
where $\bar\ell_{t+1}$ is the mean negative log-likelihood of the new task and
the small isotropic term $\tfrac{\alpha}{2N_{t+1}}\|\theta-\theta_t\|^2$ is
omitted. This is the online EWC
objective~\cite{kirkpatrick2017ewc,schwarz2018progress} with penalty strength
fixed at $N_j/N_{t+1}$, so the EWC solution is the most probable parameter under
the target of PSB with a single particle.

\subsection{Annealed Langevin transport}
\label{sec:langevin}

As $\beta$ increases from $\beta_{m-1}$ to $\beta_m$, the particle weights are
multiplied by the change in the tempered likelihood,
\begin{equation}
    w^{(k)}\leftarrow w^{(k)}\,
    p\big(\mathcal{D}_{t+1}\mid\theta^{(k)}\big)^{\beta_m-\beta_{m-1}},
    \label{eq:reweight}
\end{equation}
with each increment as large as possible while the effective number of
particles $1/\sum_k(\bar w^{(k)})^2$ of the normalized weights stays above a
fixed fraction of $K$. When the weights become uneven, low-weight particles are
replaced by copies of high-weight ones and the weights are reset. The particles
then take preconditioned Langevin steps toward $\pi_{t,\beta_m}$,
\begin{equation}
    \theta'=\theta+\eta\,A_t^{-1}\nabla_\theta\log\pi_{t,\beta_m}(\theta)
    +\sqrt{2\eta}\,A_t^{-1/2}\xi,
    \label{eq:mala}
\end{equation}
with $\xi\sim\mathcal{N}(0,I)$ and an adapted step size $\eta$, each accepted
with a probability computed from the full task likelihood. The drift of these
dynamics is the score (\ref{eq:pathscore}), and the preconditioner makes the
steps cautious in directions constrained by earlier tasks. At $\beta=1$, the
particles form the terminal population $\{\Theta_1^{(k)}\}$.

\subsection{Fisher-metric Schr\"odinger bridge}
\label{sec:bridge}

Given two marginals and a reference process, the Schr\"odinger bridge (SB) is
the path measure closest to the reference in Kullback--Leibler (KL) divergence
that connects the marginals~\cite{leonard2014survey}. At each boundary, PSB
solves
\begin{equation}
    \mathbb{P}^\star_t=\mathop{\arg\min}_{\substack{\mathbb{P}:\ \mathbb{P}_0=\hat q_t\\
    \mathbb{P}_1=\tilde q_{t+1}}}\ \mathrm{KL}(\mathbb{P}\,\|\,\mathbb{Q}_t)
    \label{eq:psb}
\end{equation}
for the Fisher-metric Brownian reference
\begin{equation}
    \mathbb{Q}_t:\quad d\Theta_s=\sigma A_t^{-1/2}\,dW_s,\qquad s\in[0,1].
    \label{eq:fisherreference}
\end{equation}
For a process with added drift $u_s$, the KL divergence to $\mathbb{Q}_t$ is
$(2\sigma^2)^{-1}\mathbb{E}\!\int_0^1\|u_s\|_{A_t}^2ds$ plus a term fixed by the
initial marginal, so moving weights that earlier tasks constrain strongly is
costly. Solving (\ref{eq:psb}) at every boundary yields a chain of SBs
$\hat q_1\to\cdots\to\hat q_T$.

We solve (\ref{eq:psb}) by iterative Markovian fitting (IMF)~\cite{shi2023dsbm}.
Given a coupling of samples $(\Theta_0,\Theta_1)$ from $\hat q_t$ and the
terminal population, intermediate points are drawn from the conditional bridge
of $\mathbb{Q}_t$,
\begin{equation}
    \Theta_s=(1-s)\,\Theta_0+s\,\Theta_1+\sigma\sqrt{s(1-s)}\,A_t^{-1/2}\xi,
    \label{eq:conditionalbridge}
\end{equation}
and a forward drift is fitted by
\begin{equation}
    \min_\phi\ \mathbb{E}\left\|v^{\mathrm{f}}_\phi(\Theta_s,s)
    -\frac{\Theta_1-\Theta_s}{1-s}\right\|_{A_t}^{2},
    \label{eq:lossf}
\end{equation}
with $s$ uniform on $[0,1]$. A backward drift $v^{\mathrm{b}}_\psi$ is fitted
analogously with target $(\Theta_0-\Theta_s)/s$. New couplings are generated by alternately
simulating the fitted forward and backward processes, and since learned drifts
only approximate the IMF projections, we report the energy distance between
simulated and terminal populations.

A drift over all weights at once would be too large to fit from $K$ particles,
so the drifts act on neuron tokens. The token of unit $i$ in a fully connected
layer $\ell$ is its incoming weights and bias,
$[W^{(\ell)}_{i,:},\,b^{(\ell)}_i]$, and in a convolutional layer it is one
filter with its bias. Any vector $u$ over the weights is split in the same way,
with $[u]^{(\ell)}_i$ its token at unit $i$ of layer $\ell$. The forward drift is
predicted token by token,
\begin{equation}
    \big[v^{\mathrm{f}}_\phi(\Theta_s,s)\big]^{(\ell)}_i
    =f^{(\ell)}_\phi\!\left([\Theta_s]^{(\ell)}_i,\,s,\,[g_s]^{(\ell)}_i,\,
    [a_s]^{(\ell)}_i\right),
    \label{eq:tokendrift}
\end{equation}
where $g_s=\nabla_\theta\log p(\mathcal{D}_{t+1}\mid\Theta_s)$, $a_s$ is the
prior term of (\ref{eq:pathscore}) at $\Theta_s$, and $f^{(\ell)}_\phi$ is a
multilayer perceptron shared by the $n_\ell$ units of layer $\ell$ that receives
$s$ through Fourier features. Sharing turns each particle into $n_\ell$ training
examples. The score inputs tell the network how the new task and the carried
posterior pull on each unit, so the drift depends on this evidence rather than
on a task identifier and remains usable at the next boundary. A zero-initialized
output layer makes the untrained drift zero, so fitting starts from the
reference (\ref{eq:fisherreference}). The backward drift has the same form.

\subsection{Bridge-guided transitions}
\label{sec:transition}

The drift fitted at one boundary guides the transport at the next. The particles
are moved by the Euler--Maruyama discretization of the previous forward drift,
with $s_m=\beta_m$ and $\Delta s_m=s_m-s_{m-1}$,
\begin{equation}
    \Theta_{m}=\Theta_{m-1}+v^{\mathrm{f}}_\phi(\Theta_{m-1},s_{m-1})\,\Delta s_m
    +\sigma\sqrt{\Delta s_m}\,A_t^{-1/2}\xi,
    \label{eq:euler}
\end{equation}
and their weights are updated by
\begin{equation}
    w^{(k)}\leftarrow w^{(k)}\,
    \frac{\pi_{t,\beta_m}(\Theta_m)\,L_m(\Theta_{m-1}\mid\Theta_m)}
         {\pi_{t,\beta_{m-1}}(\Theta_{m-1})\,K_m(\Theta_m\mid\Theta_{m-1})},
    \label{eq:bridgeweight}
\end{equation}
where $K_m$ and $L_m$ are the Gaussian transition densities of (\ref{eq:euler})
and of the corresponding backward step~\cite{delmoral2006smc}. The terminal
population thus targets $\tilde q_{t+1}$ for any drift, while an accurate drift
keeps the weights nearly uniform and reduces the Langevin steps (\ref{eq:mala})
needed after each reset. The drift is zero at the first boundary and is refitted
by IMF after every transport. The terminal particles and the updated precision
(\ref{eq:precision}) form $\hat q_{t+1}$, predictions average the particles'
predictive distributions, and the stored state, $K$ particles, one diagonal
precision and the drift networks, does not grow with the number of tasks.
.

\section{Background}
\label{sec:background}

Consider a sequence of $T$ disjoint tasks with datasets
$\mathcal{D}_1,\ldots,\mathcal{D}_T$ of i.i.d.\ samples and a model with
parameters $\theta\in\mathbb{R}^P$ shared across tasks, where the data of a task
are no longer available once it has been learned. Let $q_0(\theta)=p(\theta)$ be
the prior and $q_t(\theta)=p(\theta\mid\mathcal{D}_{1:t})$ the posterior after
task $t$. If the datasets are conditionally independent given $\theta$,
\begin{equation}
    q_{t+1}(\theta)\propto p(\mathcal{D}_{t+1}\mid\theta)\,q_t(\theta),
    \label{eq:sequential_bayes}
\end{equation}
so $q_t$ carries all the information from earlier tasks that the next update
requires. Since $q_t$ is intractable for neural networks, a replay-free learner
retains an approximation $\hat q_t$ instead. The update then defines a target
proportional to $p(\mathcal{D}_{t+1}\mid\theta)\,\hat q_t(\theta)$, which does not
directly provide samples, and errors in $\hat q_t$ propagate to all subsequent
updates~\cite{kessler2023sequential}. A replay-free learner must therefore both
represent the accumulated posterior and update this representation as new tasks
arrive.

\section{Posterior Schr\"odinger Bridges}
\label{sec:method}

PSB represents the posterior by an ensemble of networks, transports it to each
updated posterior by annealed Langevin dynamics, and fits a Schr\"odinger bridge
between consecutive ensembles that guides the next transport.

\subsection{Posterior representation and target}
\label{sec:representation}

We approximate the posterior after task $t$ by
\begin{equation}
    \hat q_t(\theta)=\sum_{k=1}^{K}w_t^{(k)}\,
    \mathcal{N}\!\left(\theta;\,\theta_t^{(k)},\,A_t^{-1}\right).
    \label{eq:mixture}
\end{equation}
Each term, which we call a component, is a Gaussian centered on a network
$\theta_t^{(k)}$, called a particle, with weight $w_t^{(k)}$. The components
cover distinct parameter solutions, and the shared diagonal precision
\begin{equation}
    A_t=\alpha I+\textstyle\sum_{j=1}^{t}\gamma^{\,t-j}N_jF_j
    \label{eq:precision}
\end{equation}
accumulates the Fisher information of each task, where $N_j=|\mathcal{D}_j|$ and
\begin{equation}
    F_j=\mathrm{diag}\Big(\mathbb{E}_{(x,y)\in\mathcal{D}_j}\mathbb{E}_{k}
    \big[\big(\nabla_\theta\log p(y\mid x,\theta_j^{(k)})\big)^{\odot2}\big]\Big)
    \label{eq:fisher}
\end{equation}
is the diagonal empirical Fisher of task $j$, with $\odot$ the elementwise
product. Here $\gamma\in(0,1]$ discounts older tasks and $\alpha>0$ is the
precision of a Gaussian prior. Since $N_jF_j$ is the Laplace precision of the
likelihood of task $j$, $A_t$ constrains each weight according to its importance
for earlier tasks. For the first task, the particles are He-uniform
initialized~\cite{he2015delving}, trained independently, and equally weighted,
treating the ensemble as approximate posterior
samples~\cite{lakshminarayanan2017ensembles,wilson2020bayesian}.

When task $t+1$ arrives, we introduce its likelihood through the tempered
posteriors
\begin{equation}
    \pi_{t,\beta}(\theta)\propto p(\mathcal{D}_{t+1}\mid\theta)^{\beta}\,
    \hat q_t(\theta),\qquad \beta\in[0,1].
    \label{eq:tempered}
\end{equation}
At $\beta=1$, this is Bayes' rule (\ref{eq:sequential_bayes}) applied to
$\hat q_t$, which defines the target $\tilde q_{t+1}=\pi_{t,1}$. Raising $\beta$
gradually avoids concentrating the weights on a few particles and moving the
particles over the full distance at once. The score of the tempered posteriors is
\begin{equation}
\begin{split}
    \nabla_\theta\log\pi_{t,\beta}(\theta)
    &= \beta\,\nabla_\theta\log p(\mathcal{D}_{t+1}\mid\theta)\\
    &\quad -A_t\Big(\theta-\textstyle\sum_{k=1}^{K}r_k(\theta)\,\theta_t^{(k)}\Big),
\end{split}
\label{eq:pathscore}
\end{equation}
where $r_k(\theta)\propto w_t^{(k)}\mathcal{N}(\theta;\theta_t^{(k)},A_t^{-1})$
is the probability that $\theta$ belongs to component $k$. The likelihood score
drives the particles toward the new task, and the prior score pulls each particle
toward the carried solutions it most likely belongs to, with a strength set by
$A_t$.

For a single component, the log-target is
$\log p(\mathcal{D}_{t+1}\mid\theta)-\tfrac12\|\theta-\theta_t\|_{A_t}^2$ up to a
constant, with $\|u\|_M^2=u^{\top}Mu$. Substituting (\ref{eq:precision}),
negating and dividing by $N_{t+1}$, maximizing the log-target is equivalent to
minimizing
\begin{equation}
    \bar\ell_{t+1}(\theta)
    +\tfrac{1}{2}\textstyle\sum_{j=1}^{t}\gamma^{\,t-j}\tfrac{N_j}{N_{t+1}}
    \|\theta-\theta_t\|_{F_j}^{2},
    \label{eq:ewc}
\end{equation}
where $\bar\ell_{t+1}$ is the mean negative log-likelihood of the new task and
the small isotropic term $\tfrac{\alpha}{2N_{t+1}}\|\theta-\theta_t\|^2$ is
omitted. This is the online EWC
objective~\cite{kirkpatrick2017ewc,schwarz2018progress} with penalty strength
fixed at $N_j/N_{t+1}$, so the EWC solution is the most probable parameter under
the target of PSB with a single particle.

\subsection{Annealed Langevin transport}
\label{sec:langevin}

As $\beta$ increases from $\beta_{m-1}$ to $\beta_m$, the particle weights are
multiplied by the change in the tempered likelihood,
\begin{equation}
    w^{(k)}\leftarrow w^{(k)}\,
    p\big(\mathcal{D}_{t+1}\mid\theta^{(k)}\big)^{\beta_m-\beta_{m-1}},
    \label{eq:reweight}
\end{equation}
with each increment as large as possible while the effective number of
particles $1/\sum_k(\bar w^{(k)})^2$ of the normalized weights stays above a
fixed fraction of $K$. When the weights become uneven, low-weight particles are
replaced by copies of high-weight ones and the weights are reset. The particles
then take preconditioned Langevin steps toward $\pi_{t,\beta_m}$,
\begin{equation}
    \theta'=\theta+\eta\,A_t^{-1}\nabla_\theta\log\pi_{t,\beta_m}(\theta)
    +\sqrt{2\eta}\,A_t^{-1/2}\xi,
    \label{eq:mala}
\end{equation}
with $\xi\sim\mathcal{N}(0,I)$ and an adapted step size $\eta$, each accepted
with a probability computed from the full task likelihood. The drift of these
dynamics is the score (\ref{eq:pathscore}), and the preconditioner makes the
steps cautious in directions constrained by earlier tasks. At $\beta=1$, the
particles form the terminal population $\{\Theta_1^{(k)}\}$.

\subsection{Fisher-metric Schr\"odinger bridge}
\label{sec:bridge}

Given two marginals and a reference process, the Schr\"odinger bridge (SB) is
the path measure closest to the reference in Kullback--Leibler (KL) divergence
that connects the marginals~\cite{leonard2014survey}. At each boundary, PSB
solves
\begin{equation}
    \mathbb{P}^\star_t=\mathop{\arg\min}_{\substack{\mathbb{P}:\ \mathbb{P}_0=\hat q_t\\
    \mathbb{P}_1=\tilde q_{t+1}}}\ \mathrm{KL}(\mathbb{P}\,\|\,\mathbb{Q}_t)
    \label{eq:psb}
\end{equation}
for the Fisher-metric Brownian reference
\begin{equation}
    \mathbb{Q}_t:\quad d\Theta_s=\sigma A_t^{-1/2}\,dW_s,\qquad s\in[0,1].
    \label{eq:fisherreference}
\end{equation}
For a process with added drift $u_s$, the KL divergence to $\mathbb{Q}_t$ is
$(2\sigma^2)^{-1}\mathbb{E}\!\int_0^1\|u_s\|_{A_t}^2ds$ plus a term fixed by the
initial marginal, so moving weights that earlier tasks constrain strongly is
costly. Solving (\ref{eq:psb}) at every boundary yields a chain of SBs
$\hat q_1\to\cdots\to\hat q_T$.

We solve (\ref{eq:psb}) by iterative Markovian fitting (IMF)~\cite{shi2023dsbm}.
Given a coupling of samples $(\Theta_0,\Theta_1)$ from $\hat q_t$ and the
terminal population, intermediate points are drawn from the conditional bridge
of $\mathbb{Q}_t$,
\begin{equation}
    \Theta_s=(1-s)\,\Theta_0+s\,\Theta_1+\sigma\sqrt{s(1-s)}\,A_t^{-1/2}\xi,
    \label{eq:conditionalbridge}
\end{equation}
and a forward drift is fitted by
\begin{equation}
    \min_\phi\ \mathbb{E}\left\|v^{\mathrm{f}}_\phi(\Theta_s,s)
    -\frac{\Theta_1-\Theta_s}{1-s}\right\|_{A_t}^{2},
    \label{eq:lossf}
\end{equation}
with $s$ uniform on $[0,1]$. A backward drift $v^{\mathrm{b}}_\psi$ is fitted
analogously with target $(\Theta_0-\Theta_s)/s$. New couplings are generated by alternately
simulating the fitted forward and backward processes, and since learned drifts
only approximate the IMF projections, we report the energy distance between
simulated and terminal populations.

A drift over all weights at once would be too large to fit from $K$ particles,
so the drifts act on neuron tokens. The token of unit $i$ in a fully connected
layer $\ell$ is its incoming weights and bias,
$[W^{(\ell)}_{i,:},\,b^{(\ell)}_i]$, and in a convolutional layer it is one
filter with its bias. Any vector $u$ over the weights is split in the same way,
with $[u]^{(\ell)}_i$ its token at unit $i$ of layer $\ell$. The forward drift is
predicted token by token,
\begin{equation}
    \big[v^{\mathrm{f}}_\phi(\Theta_s,s)\big]^{(\ell)}_i
    =f^{(\ell)}_\phi\!\left([\Theta_s]^{(\ell)}_i,\,s,\,[g_s]^{(\ell)}_i,\,
    [a_s]^{(\ell)}_i\right),
    \label{eq:tokendrift}
\end{equation}
where $g_s=\nabla_\theta\log p(\mathcal{D}_{t+1}\mid\Theta_s)$, $a_s$ is the
prior term of (\ref{eq:pathscore}) at $\Theta_s$, and $f^{(\ell)}_\phi$ is a
multilayer perceptron shared by the $n_\ell$ units of layer $\ell$ that receives
$s$ through Fourier features. Sharing turns each particle into $n_\ell$ training
examples. The score inputs tell the network how the new task and the carried
posterior pull on each unit, so the drift depends on this evidence rather than
on a task identifier and remains usable at the next boundary. A zero-initialized
output layer makes the untrained drift zero, so fitting starts from the
reference (\ref{eq:fisherreference}). The backward drift has the same form.

\subsection{Bridge-guided transitions}
\label{sec:transition}

The drift fitted at one boundary guides the transport at the next. The particles
are moved by the Euler--Maruyama discretization of the previous forward drift,
with $s_m=\beta_m$ and $\Delta s_m=s_m-s_{m-1}$,
\begin{equation}
    \Theta_{m}=\Theta_{m-1}+v^{\mathrm{f}}_\phi(\Theta_{m-1},s_{m-1})\,\Delta s_m
    +\sigma\sqrt{\Delta s_m}\,A_t^{-1/2}\xi,
    \label{eq:euler}
\end{equation}
and their weights are updated by
\begin{equation}
    w^{(k)}\leftarrow w^{(k)}\,
    \frac{\pi_{t,\beta_m}(\Theta_m)\,L_m(\Theta_{m-1}\mid\Theta_m)}
         {\pi_{t,\beta_{m-1}}(\Theta_{m-1})\,K_m(\Theta_m\mid\Theta_{m-1})},
    \label{eq:bridgeweight}
\end{equation}
where $K_m$ and $L_m$ are the Gaussian transition densities of (\ref{eq:euler})
and of the corresponding backward step~\cite{delmoral2006smc}. The terminal
population thus targets $\tilde q_{t+1}$ for any drift, while an accurate drift
keeps the weights nearly uniform and reduces the Langevin steps (\ref{eq:mala})
needed after each reset. The drift is zero at the first boundary and is refitted
by IMF after every transport. The terminal particles and the updated precision
(\ref{eq:precision}) form $\hat q_{t+1}$, predictions average the particles'
predictive distributions, and the stored state, $K$ particles, one diagonal
precision and the drift networks, does not grow with the number of tasks.
.

\section{Background}
\label{sec:background}

Consider a sequence of $T$ disjoint tasks with datasets
$\mathcal{D}_1,\ldots,\mathcal{D}_T$ of i.i.d.\ samples and a model with
parameters $\theta\in\mathbb{R}^P$ shared across tasks, where the data of a task
are no longer available once it has been learned. Let $q_0(\theta)=p(\theta)$ be
the prior and $q_t(\theta)=p(\theta\mid\mathcal{D}_{1:t})$ the posterior after
task $t$. If the datasets are conditionally independent given $\theta$,
\begin{equation}
    q_{t+1}(\theta)\propto p(\mathcal{D}_{t+1}\mid\theta)\,q_t(\theta),
    \label{eq:sequential_bayes}
\end{equation}
so $q_t$ carries all the information from earlier tasks that the next update
requires. Since $q_t$ is intractable for neural networks, a replay-free learner
retains an approximation $\hat q_t$ instead. The update then defines a target
proportional to $p(\mathcal{D}_{t+1}\mid\theta)\,\hat q_t(\theta)$, which does not
directly provide samples, and errors in $\hat q_t$ propagate to all subsequent
updates~\cite{kessler2023sequential}. A replay-free learner must therefore both
represent the accumulated posterior and update this representation as new tasks
arrive.

\section{Posterior Schr\"odinger Bridges}
\label{sec:method}

PSB represents the posterior by an ensemble of networks, transports it to each
updated posterior by annealed Langevin dynamics, and fits a Schr\"odinger bridge
between consecutive ensembles that guides the next transport.

\subsection{Posterior representation and target}
\label{sec:representation}

We approximate the posterior after task $t$ by
\begin{equation}
    \hat q_t(\theta)=\sum_{k=1}^{K}w_t^{(k)}\,
    \mathcal{N}\!\left(\theta;\,\theta_t^{(k)},\,A_t^{-1}\right).
    \label{eq:mixture}
\end{equation}
Each term, which we call a component, is a Gaussian centered on a network
$\theta_t^{(k)}$, called a particle, with weight $w_t^{(k)}$. The components
cover distinct parameter solutions, and the shared diagonal precision
\begin{equation}
    A_t=\alpha I+\textstyle\sum_{j=1}^{t}\gamma^{\,t-j}N_jF_j
    \label{eq:precision}
\end{equation}
accumulates the Fisher information of each task, where $N_j=|\mathcal{D}_j|$ and
\begin{equation}
    F_j=\mathrm{diag}\Big(\mathbb{E}_{(x,y)\in\mathcal{D}_j}\mathbb{E}_{k}
    \big[\big(\nabla_\theta\log p(y\mid x,\theta_j^{(k)})\big)^{\odot2}\big]\Big)
    \label{eq:fisher}
\end{equation}
is the diagonal empirical Fisher of task $j$, with $\odot$ the elementwise
product. Here $\gamma\in(0,1]$ discounts older tasks and $\alpha>0$ is the
precision of a Gaussian prior. Since $N_jF_j$ is the Laplace precision of the
likelihood of task $j$, $A_t$ constrains each weight according to its importance
for earlier tasks. For the first task, the particles are He-uniform
initialized~\cite{he2015delving}, trained independently, and equally weighted,
treating the ensemble as approximate posterior
samples~\cite{lakshminarayanan2017ensembles,wilson2020bayesian}.

When task $t+1$ arrives, we introduce its likelihood through the tempered
posteriors
\begin{equation}
    \pi_{t,\beta}(\theta)\propto p(\mathcal{D}_{t+1}\mid\theta)^{\beta}\,
    \hat q_t(\theta),\qquad \beta\in[0,1].
    \label{eq:tempered}
\end{equation}
At $\beta=1$, this is Bayes' rule (\ref{eq:sequential_bayes}) applied to
$\hat q_t$, which defines the target $\tilde q_{t+1}=\pi_{t,1}$. Raising $\beta$
gradually avoids concentrating the weights on a few particles and moving the
particles over the full distance at once. The score of the tempered posteriors is
\begin{equation}
\begin{split}
    \nabla_\theta\log\pi_{t,\beta}(\theta)
    &= \beta\,\nabla_\theta\log p(\mathcal{D}_{t+1}\mid\theta)\\
    &\quad -A_t\Big(\theta-\textstyle\sum_{k=1}^{K}r_k(\theta)\,\theta_t^{(k)}\Big),
\end{split}
\label{eq:pathscore}
\end{equation}
where $r_k(\theta)\propto w_t^{(k)}\mathcal{N}(\theta;\theta_t^{(k)},A_t^{-1})$
is the probability that $\theta$ belongs to component $k$. The likelihood score
drives the particles toward the new task, and the prior score pulls each particle
toward the carried solutions it most likely belongs to, with a strength set by
$A_t$.

For a single component, the log-target is
$\log p(\mathcal{D}_{t+1}\mid\theta)-\tfrac12\|\theta-\theta_t\|_{A_t}^2$ up to a
constant, with $\|u\|_M^2=u^{\top}Mu$. Substituting (\ref{eq:precision}),
negating and dividing by $N_{t+1}$, maximizing the log-target is equivalent to
minimizing
\begin{equation}
    \bar\ell_{t+1}(\theta)
    +\tfrac{1}{2}\textstyle\sum_{j=1}^{t}\gamma^{\,t-j}\tfrac{N_j}{N_{t+1}}
    \|\theta-\theta_t\|_{F_j}^{2},
    \label{eq:ewc}
\end{equation}
where $\bar\ell_{t+1}$ is the mean negative log-likelihood of the new task and
the small isotropic term $\tfrac{\alpha}{2N_{t+1}}\|\theta-\theta_t\|^2$ is
omitted. This is the online EWC
objective~\cite{kirkpatrick2017ewc,schwarz2018progress} with penalty strength
fixed at $N_j/N_{t+1}$, so the EWC solution is the most probable parameter under
the target of PSB with a single particle.

\subsection{Annealed Langevin transport}
\label{sec:langevin}

As $\beta$ increases from $\beta_{m-1}$ to $\beta_m$, the particle weights are
multiplied by the change in the tempered likelihood,
\begin{equation}
    w^{(k)}\leftarrow w^{(k)}\,
    p\big(\mathcal{D}_{t+1}\mid\theta^{(k)}\big)^{\beta_m-\beta_{m-1}},
    \label{eq:reweight}
\end{equation}
with each increment as large as possible while the effective number of
particles $1/\sum_k(\bar w^{(k)})^2$ of the normalized weights stays above a
fixed fraction of $K$. When the weights become uneven, low-weight particles are
replaced by copies of high-weight ones and the weights are reset. The particles
then take preconditioned Langevin steps toward $\pi_{t,\beta_m}$,
\begin{equation}
    \theta'=\theta+\eta\,A_t^{-1}\nabla_\theta\log\pi_{t,\beta_m}(\theta)
    +\sqrt{2\eta}\,A_t^{-1/2}\xi,
    \label{eq:mala}
\end{equation}
with $\xi\sim\mathcal{N}(0,I)$ and an adapted step size $\eta$, each accepted
with a probability computed from the full task likelihood. The drift of these
dynamics is the score (\ref{eq:pathscore}), and the preconditioner makes the
steps cautious in directions constrained by earlier tasks. At $\beta=1$, the
particles form the terminal population $\{\Theta_1^{(k)}\}$.

\subsection{Fisher-metric Schr\"odinger bridge}
\label{sec:bridge}

Given two marginals and a reference process, the Schr\"odinger bridge (SB) is
the path measure closest to the reference in Kullback--Leibler (KL) divergence
that connects the marginals~\cite{leonard2014survey}. At each boundary, PSB
solves
\begin{equation}
    \mathbb{P}^\star_t=\mathop{\arg\min}_{\substack{\mathbb{P}:\ \mathbb{P}_0=\hat q_t\\
    \mathbb{P}_1=\tilde q_{t+1}}}\ \mathrm{KL}(\mathbb{P}\,\|\,\mathbb{Q}_t)
    \label{eq:psb}
\end{equation}
for the Fisher-metric Brownian reference
\begin{equation}
    \mathbb{Q}_t:\quad d\Theta_s=\sigma A_t^{-1/2}\,dW_s,\qquad s\in[0,1].
    \label{eq:fisherreference}
\end{equation}
For a process with added drift $u_s$, the KL divergence to $\mathbb{Q}_t$ is
$(2\sigma^2)^{-1}\mathbb{E}\!\int_0^1\|u_s\|_{A_t}^2ds$ plus a term fixed by the
initial marginal, so moving weights that earlier tasks constrain strongly is
costly. Solving (\ref{eq:psb}) at every boundary yields a chain of SBs
$\hat q_1\to\cdots\to\hat q_T$.

We solve (\ref{eq:psb}) by iterative Markovian fitting (IMF)~\cite{shi2023dsbm}.
Given a coupling of samples $(\Theta_0,\Theta_1)$ from $\hat q_t$ and the
terminal population, intermediate points are drawn from the conditional bridge
of $\mathbb{Q}_t$,
\begin{equation}
    \Theta_s=(1-s)\,\Theta_0+s\,\Theta_1+\sigma\sqrt{s(1-s)}\,A_t^{-1/2}\xi,
    \label{eq:conditionalbridge}
\end{equation}
and a forward drift is fitted by
\begin{equation}
    \min_\phi\ \mathbb{E}\left\|v^{\mathrm{f}}_\phi(\Theta_s,s)
    -\frac{\Theta_1-\Theta_s}{1-s}\right\|_{A_t}^{2},
    \label{eq:lossf}
\end{equation}
with $s$ uniform on $[0,1]$. A backward drift $v^{\mathrm{b}}_\psi$ is fitted
analogously with target $(\Theta_0-\Theta_s)/s$. New couplings are generated by alternately
simulating the fitted forward and backward processes, and since learned drifts
only approximate the IMF projections, we report the energy distance between
simulated and terminal populations.

A drift over all weights at once would be too large to fit from $K$ particles,
so the drifts act on neuron tokens. The token of unit $i$ in a fully connected
layer $\ell$ is its incoming weights and bias,
$[W^{(\ell)}_{i,:},\,b^{(\ell)}_i]$, and in a convolutional layer it is one
filter with its bias. Any vector $u$ over the weights is split in the same way,
with $[u]^{(\ell)}_i$ its token at unit $i$ of layer $\ell$. The forward drift is
predicted token by token,
\begin{equation}
    \big[v^{\mathrm{f}}_\phi(\Theta_s,s)\big]^{(\ell)}_i
    =f^{(\ell)}_\phi\!\left([\Theta_s]^{(\ell)}_i,\,s,\,[g_s]^{(\ell)}_i,\,
    [a_s]^{(\ell)}_i\right),
    \label{eq:tokendrift}
\end{equation}
where $g_s=\nabla_\theta\log p(\mathcal{D}_{t+1}\mid\Theta_s)$, $a_s$ is the
prior term of (\ref{eq:pathscore}) at $\Theta_s$, and $f^{(\ell)}_\phi$ is a
multilayer perceptron shared by the $n_\ell$ units of layer $\ell$ that receives
$s$ through Fourier features. Sharing turns each particle into $n_\ell$ training
examples. The score inputs tell the network how the new task and the carried
posterior pull on each unit, so the drift depends on this evidence rather than
on a task identifier and remains usable at the next boundary. A zero-initialized
output layer makes the untrained drift zero, so fitting starts from the
reference (\ref{eq:fisherreference}). The backward drift has the same form.

\subsection{Bridge-guided transitions}
\label{sec:transition}

The drift fitted at one boundary guides the transport at the next. The particles
are moved by the Euler--Maruyama discretization of the previous forward drift,
with $s_m=\beta_m$ and $\Delta s_m=s_m-s_{m-1}$,
\begin{equation}
    \Theta_{m}=\Theta_{m-1}+v^{\mathrm{f}}_\phi(\Theta_{m-1},s_{m-1})\,\Delta s_m
    +\sigma\sqrt{\Delta s_m}\,A_t^{-1/2}\xi,
    \label{eq:euler}
\end{equation}
and their weights are updated by
\begin{equation}
    w^{(k)}\leftarrow w^{(k)}\,
    \frac{\pi_{t,\beta_m}(\Theta_m)\,L_m(\Theta_{m-1}\mid\Theta_m)}
         {\pi_{t,\beta_{m-1}}(\Theta_{m-1})\,K_m(\Theta_m\mid\Theta_{m-1})},
    \label{eq:bridgeweight}
\end{equation}
where $K_m$ and $L_m$ are the Gaussian transition densities of (\ref{eq:euler})
and of the corresponding backward step~\cite{delmoral2006smc}. The terminal
population thus targets $\tilde q_{t+1}$ for any drift, while an accurate drift
keeps the weights nearly uniform and reduces the Langevin steps (\ref{eq:mala})
needed after each reset. The drift is zero at the first boundary and is refitted
by IMF after every transport. The terminal particles and the updated precision
(\ref{eq:precision}) form $\hat q_{t+1}$, predictions average the particles'
predictive distributions, and the stored state, $K$ particles, one diagonal
precision and the drift networks, does not grow with the number of tasks.
